\DocumentMetadata{
  pdfversion=2.0,
  pdfstandard=ua-2,
  lang=en-US,
  tagging=on,
  tagging-setup={math/setup=mathml-SE,tabsorder=structure}
}
\documentclass[11pt,letterpaper]{article}
\usepackage[margin=0.68in,includefoot]{geometry}
\usepackage{newcomputermodern}
\usepackage{amsmath,amscd,mathtools}
\usepackage{tikz,adjustbox}
\providecommand{\square}{\mdlgwhtsquare}
\usepackage[hidelinks]{hyperref}
\hypersetup{
  pdftitle={Combinatorics PhD Exam, May 2022},
  pdfauthor={University of Florida Department of Mathematics},
  pdfsubject={Graduate mathematics examination},
  pdfdisplaydoctitle=true,
  bookmarksopen=true
}
\setcounter{secnumdepth}{0}
\setlength{\parindent}{0pt}
\setlength{\parskip}{0.28em}
\setlength{\emergencystretch}{3em}
\setlength{\leftmargini}{2.15em}
\setlength{\leftmarginii}{2.65em}
\setlength{\leftmarginiii}{3em}
\renewcommand{\labelenumi}{\textbf{\theenumi.}}
\pagestyle{empty}
\raggedbottom
\allowdisplaybreaks
\newenvironment{examproblems}[1]{%
  \begin{enumerate}%
  \setcounter{enumi}{#1}%
  \setlength{\topsep}{0.35em}%
  \setlength{\partopsep}{0pt}%
  \setlength{\itemsep}{0.48em}%
  \setlength{\parsep}{0.18em}%
}{\end{enumerate}}
\newenvironment{examsubparts}{%
  \begin{enumerate}%
  \setlength{\topsep}{0.18em}%
  \setlength{\partopsep}{0pt}%
  \setlength{\itemsep}{0.16em}%
  \setlength{\parsep}{0.1em}%
}{\end{enumerate}}
\newenvironment{examinstructions}{%
  \par\begingroup\itshape
}{%
  \par\endgroup\vspace{0.18em}
}
\makeatletter
\renewcommand\section{\@startsection{section}{1}{0pt}%
  {-1.25ex plus -.25ex minus -.1ex}%
  {0.55ex plus .1ex}%
  {\normalfont\large\bfseries}}
\renewcommand{\@maketitle}{%
  \newpage\null\vskip -1em%
  {\footnotesize\bfseries
    \href{https://www.ufl.edu/}{University of Florida}, \href{https://math.ufl.edu/}{Department of Mathematics}\par}%
  \vskip -0.5em%
  \tagstructbegin{tag=Title}%
    {\LARGE\bfseries\href{https://gma.math.ufl.edu/exams/combinatorics-phd/}{Combinatorics PhD Exam}, May 2022\par}%
  \tagstructend%
  \vskip 0.25em%
  \tagmcbegin{artifact}%
    \hrule height 1pt%
  \tagmcend%
  \vskip 1em%
}
\makeatother
\title{Combinatorics PhD Exam, May 2022}
\author{}
\date{}
\begin{document}
\maketitle
\thispagestyle{empty}
\begin{examproblems}{0}
\item
Let \(G_n\) be the graph whose vertex set is the set of all permutations of \([n]\) (in one line notation), with two vertices adjacent if they differ by switching two consecutive entries in the permutation. What is the chromatic number of \(G_n\) (with proof).
\item
Let \(T=(V,E)\) be a tournament that is not strongly connected. Show there is a partition \(V=A\sqcup B\) with \(A,B\) non-empty so that every edge between \(a\in A\) and \(b\in B\) is from~\(a\) to~\(b\).
\item
Recall the permutation \(\pi=\pi_1\cdots\pi_n\) contains the permutation \(\sigma=\sigma_1\cdots\sigma_k\) if there is a subset \(\{i_1<\cdots<i_k\}\subseteq[n]\) so that \(\pi_{i_1}\cdots\pi_{i_k}\) has the same relative order as~\(\sigma\), and that~\(\pi\) avoids~\(\sigma\) if it does not contain~\(\sigma\). Let \(S_n(\sigma)\) be the set of permutations avoiding~\(\sigma\). Prove that \(\lvert S_n(231)\rvert\) is the~\(n\)th Catalan number \(C_n=\binom{2n}{n}/(n+1)\).
\item
Determine the number of binary strings of length~\(n\) beginning with~\(0\), ending with~\(1\), and such that the number of copies of \(00\) equals the number of copies of \(11\). For example, \(00011011\) has two copies of each.
\item
Show that the number of partitions of~\(n\) with no part divisible by~\(k\) is equal to the number of partitions of~\(n\) with each part appearing at most \(k-1\) times.
\item
Let~\(P\) be a finite poset with unique minimal element~\(0\) and Möbius function~\(\mu\). Let \(y\in P\) so that there is only one~\(x\) covered by~\(y\). Show \(\mu(0,y)=0\).
\item
Let \(e_k(x_1,\ldots,x_n)\) be the elementary symmetric function in~\(n\) variables and \(c(n,k)\) be the (signless) Stirling number of the first kind. Prove that 
\[
e_{n+1-k}(1,2,\ldots,n)=c(n+1,k).
\]
 For example, \(e_2(1,2,3)=1\cdot2+1\cdot3+2\cdot3=11=c(4,2)\).

\par
(Hint: what is \(e_k(x_1,\ldots,x_n)\) in terms of symmetric functions in \(n-1\) variables?)
\item
Let \(D_n\) be the set of permutations of~\(n\) with no fixed points and \(F_n\) be the set of permutations with exactly one fixed point. Show 
\[
\lvert D_n\rvert-\lvert F_n\rvert=(-1)^n.
\]
\end{examproblems}
\end{document}
